\section{Discussion}\label{sec:discussion}
The implementation shows that a closed, traceable three-verdict contract is feasible for Python functions. The data, however, do not show that the distinguishing element, the sufficiency level, adds value. Plain repetition reproduces V7-B's flakiness handling at similar protection; V7-B's lower false assurance is largely relocation to \INS{}, and \INS{} is itself dominated by features the engine does not implement. Where V7-B differs from B1 most visibly (renamed entries) the difference is a refusal rather than a finding. An independent review's readiness verdict was that the work is publishable only as a design, feasibility or negative-result report, with the honest alternative that the contribution is better framed as an evidence contract for differential testing than as a sufficiency advance. We adopt that framing.

What a decisive study would need: independent labelled real pairs with a kappa-checked truth team; a sealed holdout; matched-coverage false-assurance comparison of V7-B against B1+rep; at least one published tool (Testora, DiffTestGen or Mokav class) or dropping competitive claims; and engine support for renames, generators and class units so that \INS{} carries meaning. Current data predict, but do not establish, that sufficiency would show little gain over B1+rep.
